\section{Introduction}
\label{sec:intro}

Diffusion models have emerged as powerful generative models for producing high-quality data samples, ranging from images to audio and beyond~\citep{ho2020denoising, song2020denoising, song2020score}. These models work by iteratively transforming a simple noise distribution into complex, structured outputs, guided by a learned reverse diffusion process. Despite their effectiveness, diffusion models come with significant computational costs~\citep{liu2022pseudo, li2023q}. The iterative nature of the sampling process, which requires multiple inferences through neural networks, makes these models computationally expensive and time-intensive. This limitation restricts their scalability and accessibility.

Existing works aim to enhance the efficiency of the sampling process in diffusion models through several strategies. Caching methods, for example, accelerate diffusion models by reusing intermediate computations~\citep{ma2024deepcache, wimbauer2024cache}. These methods exploit the significant similarities between features at nearby time steps, enabling the skipping of redundant computations by directly using cached results. Additionally, quantization techniques reduce inference costs by converting model weights and activations into integers using scaling factors~\citep{nagel2021white, yang2019quantization}. Among these, post-training quantization (PTQ) stands out as a promising approach since it estimates scaling factors in a training-free manner, making it broadly applicable to pre-trained models~\citep{li2021brecq}. Another line of work focuses on efficient sampling strategies with solvers or samplers, such as denoising diffusion implicit models (DDIMs), which significantly reduce the number of sampling steps required in diffusion models, and speed up the process~\citep{song2020denoising}.

Our preliminary studies reveal that while caching and PTQ methods have achieved notable success in accelerating the sampling process, they also introduce significant limitations. First, our analysis reveals that caching methods can lead to reuse errors that accumulate throughout the generation process, particularly when reuse schedules are not carefully designed with respect to the time step and reused components. For instance, when following the reuse strategy of DeepCache~\citep{ma2024deepcache}, but slightly modifying the reused components, we observe that the relative $\ell_2$ distance between the features of a standard diffusion model and those of caching methods increases significantly throughout the generation process, reaching 40\% in the final step, even when the cache is updated every three steps. On the other hand, our studies show that diffusion models exhibit significant outliers in activations and variations in activation ranges across time steps, leading to substantial quantization errors under 
low-bit activation quantization. 

In this work, we propose \textbf{Modulated Diffusion (MoDiff)}, an innovative, rigorous, and principled framework that accelerates the diffusion sampling process while addressing the limitations of existing methods. 
Specifically, we propose modulated computation to significantly reduce activation quantization error by leveraging the computation redundancy across the diffusion time steps. Moreover, we further introduce a novel error compensation modulation to address error accumulation. Furthermore, we provide theoretical analyses to explain why the temporal difference results in lower quantization error and how error compensation effectively eliminates accumulated errors.
Our extensive experiments validate the effectiveness of this framework, demonstrating that MoDiff pushes the activation quantization limit of PTQ methods from 8 bits to as low as 3 bits without any performance degradation, all within a training-free manner.

The proposed MoDiff framework inherits the advantages of existing acceleration methods while addressing their limitations. It significantly generalizes the caching techniques through modulated computation but reduces apprpoximation and accumulated error. Additionally, from a quantization perspective, MoDiff significantly reduces the quantization error of existing PTQ methods, enabling the use of much lower activation bit-widths without sacrificing performance. Notably, MoDiff is agnostic to quantization algorithms and can be generally applied across different methods, making its contribution orthogonal to existing PTQ techniques. Furthermore, MoDiff imposes no constraints on samplers, ensuring compatibility with solver-based acceleration methods.
In summary, our main contributions are as follows:

\begin{itemize}
\item We present an in-depth and insightful preliminary study that reveals the limitations of existing acceleration techniques for diffusion models, such as caching and quantization methods, highlighting issues like error accumulation and high approximation error. 
\item We propose MoDiff, a novel, rigorous, and principled framework that accelerates diffusion models through modulated quantization and error compensation. MoDiff not only inherits the advantages of existing methods but also overcomes their limitations, enabling significantly more aggressive activation quantization.
\item We provide theoretical analyses of quantization error and the error compensation mechanism in MoDiff, demonstrating that our approach can significantly reduce the required activation bit precision in PTQ.
\item Extensive experiments on CIFAR-10, LSUN-Churches, and LSUN-Bedroom show that MoDiff enables state-of-the-art quantization techniques to reduce activation precision from 8 bits to as low as 3 bits without any performance degradation in a training-free manner.
\end{itemize}


